% Chapter 3 -- Rubric: PO2 (2.1), PO3 (3.1, 3.2)
\chapter{Requirements Analysis}
\label{ch:requirements}

\section{Project Objectives and Constraints}
\label{sec:requirements}
% \guidance{Identify all stakeholders (users, client or industry partner,
% administrators, regulators) and their expectations, including conflicting
% ones. Then clearly state the objectives, the functional and non-functional
% requirements, and the constraints (time, budget, platform, data,
% deployment), taking those expectations into account.}

\subsection{Problem Context and Origin}
\label{subsec:problem-context}

The project originated with Spectrum Software \& Consulting Pvt. Ltd. (SSCL), a software development organisation. Engineering leadership at
SSCL monitored the health of their software projects through four to five
separate tool dashboards such as a version control platform (e.g., Github), an issue-tracking and sprint-planning platform (e.g., Jira), a static code-analysis service (e.g., SonarQube), and a CI/CD service (e.g., Jenkins). Each dashboard presented its own metrics in its
own vocabulary, with no shared notion of whether a project was, overall, in good
or poor health. Assessing the state of a single project required visiting every
dashboard in turn and mentally reconciling the results; assessing a portfolio of
projects this way was impractical.

SSCL proposed a platform that would ingest data from these separate tools,
translate heterogeneous raw metrics into a common, comparable scale, and present
the result in a single view. The system described in this report addresses that
problem. It is explicitly positioned as a trend indicator system rather
than a real-time monitoring tool. Its purpose is to reveal how projects and teams
are trending over weeks and months, not to report the state of a build in the
last five minutes.

\subsection{Requirements Elicitation}
\label{subsec:elicitation}

Requirements were gathered through three complementary channels. The primary
channel was a series of meetings with the Chief Technical Officer of
SSCL, who originated the concept and articulated the operational pain point,
the tools in use, and the decisions the system was intended to support. A
secondary channel was an intern at SSCL with prior exposure to the concept,
who provided practical grounding on how the organisation's teams actually work
day to day. The third channel was the team's own analysis such as study of the public
APIs exposed by the candidate tools to establish which metrics were obtainable in
practice, and derivation of a scoring model from that feasible metric set.
Throughout, the university supervisor provided guidance on scope, methodology and
academic rigour.

\subsection{Stakeholders and Their Expectations}
\label{subsec:stakeholders}

\begin{longtable}{
  >{\raggedright\arraybackslash}p{0.08\textwidth}
  >{\raggedright\arraybackslash}p{0.19\textwidth}
  >{\raggedright\arraybackslash}p{0.17\textwidth}
  >{\raggedright\arraybackslash}p{0.40\textwidth}
}
\caption{Stakeholders and their principal expectations}
\label{tab:stakeholders}\\
\toprule
\textbf{ID} & \textbf{Stakeholder} & \textbf{Role} & \textbf{Principal expectations} 
\\
\midrule                                                                             
\endfirsthead
\toprule                                                                             
\textbf{ID} & \textbf{Stakeholder} & \textbf{Role} & \textbf{Principal expectations} \\
\midrule
\endhead
\bottomrule
\endlastfoot
SH-01 & SSCL leadership & Client; originator; primary requirement source
& A single consolidated view replacing four to five dashboards; comparable scores
across projects; early visibility of declining projects \\
\addlinespace
SH-02 & Engineering managers and project managers & Primary end users & Status at a glance; the ability to drill into a score that looks wrong; low operational overhead \\
\addlinespace
SH-03 & Developers employed by SSCL & Data subjects and survey respondents & Assurance that the system is not an instrument of individual surveillance or ranking; minimal interruption; the ability to answer candidly without attribution \\
\addlinespace
SH-04 & Company administrator & Tenant and credential owner & Control over access tokens, workspace membership and tool configuration; assurance that credentials are not exposed \\
\addlinespace
SH-05 & University supervisor & Academic oversight & Methodological rigour, documented decision-making, demonstrable evaluation \\
\addlinespace
SH-06 & External tool and service providers & Constraint setters & Adherence to their authentication models, rate limits, free-tier quotas and terms of service \\
\addlinespace
SH-07 & Development team & Implementers & A scope deliverable by five students within two semesters\\
\end{longtable}

\subsection{Conflicting Expectations and Their Resolution}
\label{subsec:conflicts}

Several stakeholder expectations are in genuine tension. Each conflict was
resolved by an explicit requirement rather than left to implementation judgement.

\begin{longtable}{
  >{\raggedright\arraybackslash}p{0.04\textwidth}
  >{\raggedright\arraybackslash}p{0.34\textwidth}
  >{\raggedright\arraybackslash}p{0.34\textwidth}
  >{\raggedright\arraybackslash}p{0.13\textwidth}
}
\caption{Conflicting stakeholder expectations and their resolution}
\label{tab:conflicts}\\
\toprule
\textbf{\#} & \textbf{Conflict} & \textbf{Resolution} & \textbf{Encoded by} \\
\midrule
\endfirsthead
\toprule
\textbf{\#} & \textbf{Conflict} & \textbf{Resolution} & \textbf{Encoded by} \\
\midrule
\endhead
\bottomrule
\endlastfoot
1 & SH-01 seeks insight into how teams are performing; SH-03, the same
organisation's employees, require that the system not attribute opinions to
individuals & Qualitative feedback is collected anonymously, with no stored link
between a respondent and their answers, and aggregate insight is withheld below a
minimum response count & NFR-07, NFR-08 \\
\addlinespace
2 & SH-01 and SH-02 would benefit from frequent team feedback; SH-03 experience
repeated surveys as interruption and survey fatigue & Survey frequency is capped per
project per month, and a cooldown is enforced per developer across all projects, with
fair rotation for developers on several projects & FR-33, FR-31 \\
\addlinespace
3 & SH-01 expects current data; SH-06 impose rate limits that make continuous polling
infeasible under free-tier quotas & The product is positioned as a trend indicator
on a scheduled cadence rather than a real-time monitor, with on-demand sync available
when a current reading is needed & FR-13, FR-18, NFR-16 \\
\addlinespace
4 & SH-01 expects every project to carry a comparable score; in practice most
projects have only a subset of the four tool categories configured & Scores are
computed from whichever inputs are available, with the weight of absent inputs
redistributed proportionally rather than treated as zero & FR-22 \\
\addlinespace
5 & SH-04 must be able to verify and rotate access tokens; the same tokens must never
be casually exposed & Credentials are masked by default in all responses and
revealed only through an explicit, administrator-restricted action, and are encrypted
at rest & FR-12, NFR-01, NFR-04 \\
\addlinespace
6 & SH-01 expects the analytical capability that AI features provide; C-03 (in Section~\ref{subsec:constraints}) permits no
expenditure, and external AI services may be unavailable & AI-dependent features are
optional, the system degrades to non-AI behaviour when such services are
unconfigured or unavailable, rather than failing & NFR-13 \\
\addlinespace
7 & SH-01 requested coverage of their full toolchain; SH-07 capacity permits only a
subset within the available time & A representative connector from each of the four
tool categories was implemented, behind an abstraction that permits further providers
to be added without altering scoring logic & NFR-21 \\
\addlinespace
8 & SH-05 values methodological rigour and documentation; SH-01 values a working,
maintainable product & Both were treated as delivery criteria: design decisions are
documented and the scoring rules are specified and unit-tested, which serves the
academic requirement while improving maintainability for the client & NFR-22, NFR-23
\\
\end{longtable}

\subsection{Project Objectives}
\label{subsec:objectives}

\begin{longtable}{
  >{\raggedright\arraybackslash}p{0.08\textwidth}
  >{\raggedright\arraybackslash}p{0.78\textwidth}
}
\caption{Project objectives}
\label{tab:objectives}\\
\toprule
\textbf{ID} & \textbf{Objective} \\
\midrule
\endfirsthead
\toprule
\textbf{ID} & \textbf{Objective} \\
\midrule
\endhead
\bottomrule
\endlastfoot
O-01 & Consolidate engineering metrics from independent tools into a single
dashboard, removing the need to consult each tool's own dashboard for routine status
assessment \\
\addlinespace
O-02 & Normalise heterogeneous raw metrics into comparable, interpretable health
scores on a uniform 0--100 scale across seven quality dimensions \\
\addlinespace
O-03 & Keep every score explainable and auditable, so that a manager can determine
which underlying metrics produced a score and how much each contributed \\
\addlinespace
O-04 & Capture qualitative team signal via survey that quantitative metrics cannot express,
without compromising the anonymity of individual respondents \\
\addlinespace
O-05 & Close the management loop by recording the actions taken in response to
identified problems and reviewing their effectiveness after a defined interval \\
\addlinespace
O-06 & Operate unattended, accumulating trend data on a fixed cadence without
requiring manual triggering \\
\addlinespace
O-07 & Remain extensible to additional tools and providers without modification of
the core scoring logic \\
\end{longtable}

Objectives O-03 and O-05 distinguish the system from a conventional metrics
dashboard: the first makes scores accountable rather than opaque, and the second
records the managerial response alongside the measurement, enabling later
assessment of whether interventions worked.

\subsection{Functional Requirements}
\label{subsec:functional-requirements}

\subsubsection{Authentication, Tenancy and Access Control}

\begin{longtable}{
  >{\raggedright\arraybackslash}p{0.08\textwidth}
  >{\raggedright\arraybackslash}p{0.78\textwidth}
}
\caption{Functional requirements: authentication, tenancy and access control}
\label{tab:fr-auth}\\
\toprule
\textbf{ID} & \textbf{Requirement} \\
\midrule
\endfirsthead
\toprule
\textbf{ID} & \textbf{Requirement} \\
\midrule
\endhead
\bottomrule
\endlastfoot
FR-01 & The system shall allow a new organisation to register, creating a company and
an initial administrator account, with the administrator's email address verified by
a one-time code before the account is created \\
\addlinespace
FR-02 & The system shall authenticate users by email address and password and
maintain an authenticated session \\
\addlinespace
FR-03 & The system shall allow a user to reset a forgotten password through a
single-use emailed link, and shall invalidate all of that user's existing sessions
upon reset \\
\addlinespace
FR-04 & The system shall allow an administrator to invite a person to a project by
email address, using a single-use invitation that the recipient accepts either by
registering or by accepting while already authenticated \\
\addlinespace
FR-05 & The system shall support two roles, administrator and member, where a member
may access only the projects to which they are assigned and an administrator may
access all projects within their company \\
\addlinespace
FR-06 & The system shall allow an administrator to remove a member from a project \\
\end{longtable}

\subsubsection{Workspace and Project Configuration}

\begin{longtable}{
  >{\raggedright\arraybackslash}p{0.08\textwidth}
  >{\raggedright\arraybackslash}p{0.78\textwidth}
}
\caption{Functional requirements: workspace and project configuration}
\label{tab:fr-workspace}\\
\toprule
\textbf{ID} & \textbf{Requirement} \\
\midrule
\endfirsthead
\toprule
\textbf{ID} & \textbf{Requirement} \\
\midrule
\endhead
\bottomrule
\endlastfoot
FR-07 & The system shall allow an administrator to create a workspace defined by a
version-control provider, an organisation identifier and an access token \\
\addlinespace
FR-08 & The system shall list the repositories reachable with a supplied access token
before any data is stored, so the administrator can confirm the token's scope \\
\addlinespace
FR-09 & The system shall import selected repositories as projects belonging to the
workspace \\
\addlinespace
FR-10 & The system shall allow further repositories to be imported into an existing
workspace, excluding those already imported \\
\addlinespace
FR-11 & The system shall allow an administrator to configure per-project integrations
for project management, code quality and continuous integration tools, with the
version-control credential inherited from the workspace \\
\addlinespace
FR-12 & The system shall mask credential values in all responses by default, and
shall reveal a credential only through an explicit administrator-restricted request
\\
\end{longtable}

\subsubsection{Metric Ingestion}

\begin{longtable}{
  >{\raggedright\arraybackslash}p{0.08\textwidth}
  >{\raggedright\arraybackslash}p{0.78\textwidth}
}
\caption{Functional requirements: metric ingestion}
\label{tab:fr-ingestion}\\
\toprule
\textbf{ID} & \textbf{Requirement} \\
\midrule
\endfirsthead
\toprule
\textbf{ID} & \textbf{Requirement} \\
\midrule
\endhead
\bottomrule
\endlastfoot
FR-13 & The system shall allow a user to trigger an immediate synchronisation of a
project across all of its configured tools \\
\addlinespace
FR-14 & The system shall retrieve metrics from a version-control platform, an
issue-tracking platform, a static code-analysis service and a continuous-integration
service \\
\addlinespace
FR-15 & The system shall record each synchronisation as an immutable, timestamped
project snapshot with the metrics retrieved from each tool \\
\addlinespace
FR-16 & The system shall report synchronisation progress to the user interface in
real time, indicating the status of each tool individually \\
\addlinespace
FR-17 & The system shall allow a user to navigate away from or reload the page during
a synchronisation and reattach to its progress \\
\addlinespace
FR-18 & The system shall automatically synchronise every configured project on a
configurable daily schedule without user intervention \\
\addlinespace
FR-19 & The system shall complete a synchronisation and record the data successfully
retrieved when one or more individual tools fail, reporting the outcome as partial \\
\end{longtable}

\subsubsection{Health Scoring and Presentation}

\begin{longtable}{
  >{\raggedright\arraybackslash}p{0.08\textwidth}
  >{\raggedright\arraybackslash}p{0.78\textwidth}
}
\caption{Functional requirements: health scoring and presentation}
\label{tab:fr-scoring}\\
\toprule
\textbf{ID} & \textbf{Requirement} \\
\midrule
\endfirsthead
\toprule
\textbf{ID} & \textbf{Requirement} \\
\midrule
\endhead
\bottomrule
\endlastfoot
FR-20 & The system shall compute seven health scores per snapshot --- security,
reliability, maintainability, CI/CD and deployment health, team health, engineering
process, and planning and execution --- each on a 0--100 scale where a higher value
indicates better health \\
\addlinespace
FR-21 & The system shall compute an overall health score for each snapshot from the
seven dimension scores \\
\addlinespace
FR-22 & The system shall compute scores from the metrics actually available for a
snapshot, redistributing the weight of unavailable metrics across those present \\
\addlinespace
FR-23 & The system shall persist computed scores against the snapshot from which they
were derived \\
\addlinespace
FR-24 & The system shall present a portfolio view listing all accessible projects
with their overall score, dimension scores and trend \\
\addlinespace
FR-25 & The system shall present a per-project dashboard showing the current score,
its trend, the dimension scores and their history over time \\
\addlinespace
FR-26 & The system shall present, on request, a breakdown of any dimension score
identifying the contributing metrics, their measured values, their applied weights
and their individual contributions \\
\end{longtable}

\subsubsection{Team Feedback Surveys}

\begin{longtable}{
  >{\raggedright\arraybackslash}p{0.08\textwidth}
  >{\raggedright\arraybackslash}p{0.78\textwidth}
}
\caption{Functional requirements: team feedback surveys}
\label{tab:fr-surveys}\\
\toprule
\textbf{ID} & \textbf{Requirement} \\
\midrule
\endfirsthead
\toprule
\textbf{ID} & \textbf{Requirement} \\
\midrule
\endhead
\bottomrule
\endlastfoot
FR-27 & The system shall generate survey questions automatically, informed by the
project's current health scores, their trend since the previous synchronisation, andnotable operational conditions detected in the latest metrics \\
\addlinespace FR-28 & The system shall filter generated questions for duplication and assess their
quality, discarding those below a configurable quality threshold \\
\addlinespace
FR-29 & The system shall allow an authorised reviewer to edit the generated questions before dispatch, and shall prevent further editing once the survey has been distributed \\ \addlinespace                                                               FR-30 & The system shall distribute one scheduled pulse survey per project per month, at a time selected within a configurable window and fixed once assigned \\
\addlinespace FR-31 & The system shall allow an authorised user to dispatch an additional survey on
demand, subject to a configurable monthly limit per project \\
\addlinespace
FR-32 & The system shall distribute a survey as a single anonymous link, broadcast to
configured team communication channels and emailed to eligible project developers \\
\addlinespace                                                                        
FR-33 & The system shall not email a developer a survey within a configurable
cooldown period of their previous survey email, and shall rotate fairly between the  
projects a developer belongs to \\
\addlinespace                                                                        
FR-34 & The system shall accept anonymous responses through a time-limited link that
expires at a configurable deadline \\                                                
\addlinespace
FR-35 & The system shall close a survey at its deadline or on demand and produce an  
analysis comprising dimension scores, recurring themes, a narrative summary and a
per-question summary \\                                                              
\addlinespace
FR-36 & The system shall flag a project as due for a survey when defined health      scores fall below configured thresholds \\
\end{longtable}

\subsubsection{Management Action Log}

\begin{longtable}{
  >{\raggedright\arraybackslash}p{0.08\textwidth}
  >{\raggedright\arraybackslash}p{0.78\textwidth}
}
\caption{Functional requirements: management action log}
\label{tab:fr-actions}\\
\toprule
\textbf{ID} & \textbf{Requirement} \\
\midrule
\endfirsthead
\toprule
\textbf{ID} & \textbf{Requirement} \\
\midrule
\endhead
\bottomrule
\endlastfoot
FR-37 & The system shall allow a user to record a management action comprising the
problem observed, its assessed cause, the action taken, the date, and the projects
affected \\
\addlinespace
FR-38 & The system shall list recorded actions per project and across the company,
scoped to the user's role \\
\addlinespace
FR-39 & The system shall allow actions to be searched by keyword, and shall
additionally offer a semantic search that retrieves conceptually related actions \\
\addlinespace
FR-40 & The system shall prompt the author of an action to rate its effectiveness
after a defined review interval, and shall allow that review to be deferred \\
\addlinespace
FR-41 & The system shall allow a member to modify or delete actions they recorded,
and an administrator to do so for any action within their company \\
\end{longtable}

\subsection{Non-Functional Requirements}
\label{subsec:non-functional-requirements}

\begin{longtable}{
  >{\raggedright\arraybackslash}p{0.08\textwidth}
  >{\raggedright\arraybackslash}p{0.78\textwidth}
}
\caption{Non-functional requirements}
\label{tab:nfr}\\
\toprule
\textbf{ID} & \textbf{Requirement} \\
\midrule
\endfirsthead
\toprule
\textbf{ID} & \textbf{Requirement} \\
\midrule
\endhead
\bottomrule
\endlastfoot
\multicolumn{2}{l}{\textit{Security}} \\
\addlinespace
NFR-01 & Access tokens and credential material shall be encrypted at rest \\
\addlinespace
NFR-02 & User passwords shall be stored only as salted cryptographic hashes and shall
never be recoverable \\
\addlinespace
NFR-03 & Session identifiers shall be transmitted in cookies inaccessible to
client-side scripts, shall expire after a defined period of inactivity, and shall be
revocable \\
\addlinespace
NFR-04 & Credential values shall not appear in API responses except through the
explicit reveal action of FR-12, and shall never be written to application logs \\
\addlinespace
NFR-05 & Endpoints accessible without authentication shall be rate-limited \\
\addlinespace
NFR-06 & Survey access links shall be tamper-evident, such that modification of an
embedded identifier or deadline invalidates the link \\
\addlinespace
\multicolumn{2}{l}{\textit{Privacy}} \\
\addlinespace
NFR-07 & A survey response shall not be attributable to the individual who submitted
it, and no stored association between respondent identity and response content shall
exist \\
\addlinespace
NFR-08 & Aggregate survey insight shall be reported only where the number of
responses meets a configured minimum \\
\addlinespace
NFR-09 & Data transmitted to third-party analysis services shall not include the
names of individuals or identifiers of specific work items \\
\addlinespace
\multicolumn{2}{l}{\textit{Reliability}} \\
\addlinespace
NFR-10 & The failure of one tool during a synchronisation shall not prevent the
remaining tools from being processed and recorded \\
\addlinespace
NFR-11 & Background jobs shall be retried automatically on failure with an increasing
delay between attempts \\
\addlinespace
NFR-12 & Scheduled synchronisation shall be idempotent, such that a repeated or
retried schedule trigger does not duplicate work for the same project and time slot
\\
\addlinespace
NFR-13 & The system shall degrade gracefully when optional external services are
unconfigured or unavailable, substituting reduced functionality rather than failing
\\
\addlinespace
NFR-14 & Multi-step operations that cannot be performed atomically shall reverse
their partial effects on failure \\
\addlinespace
\multicolumn{2}{l}{\textit{Performance and Efficiency}} \\
\addlinespace
NFR-15 & Calls to external APIs shall be issued with bounded concurrency \\
\addlinespace
NFR-16 & The system shall observe the rate limits published by external providers,
and shall stagger projects that share a credential so that a scheduled run does not
exhaust a single credential's quota \\
\addlinespace
NFR-17 & Dashboard views shall be served from stored snapshot data and shall not
block on live calls to external tools \\
\addlinespace
\multicolumn{2}{l}{\textit{Usability}} \\
\addlinespace
NFR-18 & A manager shall be able to assess routine project status without opening any
of the underlying tools \\
\addlinespace
NFR-19 & The progress of a running synchronisation shall be visible to the user, with
per-tool status \\
\addlinespace
NFR-20 & Any presented score shall be explainable on demand in terms of its
contributing metrics \\
\addlinespace
\multicolumn{2}{l}{\textit{Maintainability}} \\
\addlinespace
NFR-21 & Support for an additional tool provider shall be addable without
modification of the scoring logic \\
\addlinespace
NFR-22 & The scoring rules for each health dimension shall be documented and covered
by automated tests \\
\addlinespace
NFR-23 & Every change shall be validated automatically for type correctness, static
analysis compliance, successful build and passing tests before it is merged \\
\addlinespace
\multicolumn{2}{l}{\textit{Portability and Deployability}} \\
\addlinespace
NFR-24 & All environment-specific configuration shall be supplied externally, with no
credential or endpoint hard-coded in source \\
\addlinespace
NFR-25 & The server components shall be deployable as containers \\
\addlinespace
NFR-26 & Scheduling shall be configurable to an arbitrary named time zone,
independent of the host's clock \\
\end{longtable}

\subsection{Constraints}
\label{subsec:constraints}

\begin{longtable}{
  >{\raggedright\arraybackslash}p{0.06\textwidth}
  >{\raggedright\arraybackslash}p{0.13\textwidth}
  >{\raggedright\arraybackslash}p{0.67\textwidth}
}
\caption{Project constraints}
\label{tab:constraints}\\
\toprule
\textbf{ID} & \textbf{Category} & \textbf{Constraint and implication} \\
\midrule
\endfirsthead
\toprule
\textbf{ID} & \textbf{Category} & \textbf{Constraint and implication} \\
\midrule
\endhead
\bottomrule
\endlastfoot
C-01 & Time & Development spanned two academic semesters, approximately twelve
months, undertaken alongside concurrent coursework rather than full-time. This
required incremental delivery and forced explicit deferral of lower-priority
capability \\
\addlinespace
C-02 & Human resources & The team comprised five students with no dedicated
operations, quality-assurance or design specialists; all such responsibilities were
absorbed by the same five members \\
\addlinespace
C-03 & Budget & No monetary budget was available. Every external service ---
database, cache, hosting, AI inference, email delivery and code analysis --- had to
operate within a free service tier. This constrained hosting to a single application
instance without redundancy, and bounded the volume of AI inference and external API
traffic the system could consume \\
\addlinespace
C-04 & Platform & The implementation targets a Node.js runtime with TypeScript.
Persistence uses a managed PostgreSQL service accessed through a REST client, which
does not expose multi-statement transactions across tables; operations spanning
several tables therefore require explicit compensating logic. An in-memory data store
is required for both job queuing and session storage, making it a hard runtime
dependency \\
\addlinespace
C-05 & Data availability & The system can only measure what the external tools
expose. Several intended metrics are unobtainable in practice: security alert counts
are unavailable where the repository has not enabled the relevant scanning feature;
test coverage is reported only where the project's pipeline publishes a coverage
report; and certain issue-tracker fields are instance-specific custom fields whose
identifiers vary between organisations \\
\addlinespace
C-06 & Data history & No historical backfill is possible from the source tools for most metrics. Trend data therefore accumulates only from the first synchronisation onward, meaning the system's analytical value increases over time and is limited immediately after adoption \\
\addlinespace
C-07 & Integration access & The system operates strictly read-only against the partner's tools using access tokens, with no ability to install agents, configure webhooks or otherwise modify the partner's infrastructure \\
\addlinespace
C-08 & Deployment & The system is deployed to managed cloud platforms --- a container platform for the server components, a static hosting platform for the client application, and managed cache and database services. At the time of writing it is hosted and operational but has not yet been handed over to [PARTNER] for production use \\
\addlinespace
C-09 & Confidentiality & Work is conducted under a non-disclosure agreement with [PARTNER], which governs the handling of their data and the disclosure of information in this report. The obligations arising from it are addressed in Section~\ref{sec:standards} \\
\end{longtable}

\section{Regulatory Requirements and Standards}
\label{sec:standards}
\guidance{Identify all relevant regulatory requirements, engineering standards, and
codes of practice (e.g. data-protection law, WCAG, OWASP, ISO/IEC/IEEE
standards, industry rules, partner policies). For each, state why it applies
and how the design complies.}

\section{Formulation of Solutions}
\label{sec:alternatives}
\guidance{Formulate multiple genuinely different solutions that could each
functionally meet most requirements. Compare and evaluate them against
explicit criteria from the requirements and constraints, and draw a valid
conclusion that justifies the selected solution.}